\subsection{Uniform finite bounds for typed evaluation}
\label{obs:bound-section}

We verify that the preceding clauses define actual values in the stated
spaces, including at free minus layers.  In this subsection a typing tree
is fixed, but its typed substitution and its parameters may vary.  Bounds
will depend only on the finite tree and an integer, and not on any of
these images or parameters.  We use the images of $C_k$ in $C_*$ without
further notation.

Two environments for the same tree are called \emph{image-compatible} if their
term images are pointwise convertible.  Their actual normalized
declaration types consequently agree.  For $r\geq 0$, say that they
\emph{agree to level $r$} if, in addition, corresponding plain and minus
parameters are equal and corresponding plus parameters are
$\equiv_r$-equivalent.  Parameters at inactive declarations are absent.
An environment is \emph{minus-bounded by $k$} if all its minus parameters
belong to $C_k$.  In particular, being minus-bounded by zero means having
no minus parameters.  Every finite environment is minus-bounded by some
$k$, since $C_*$ is the increasing union of its stages.

\paragraph{The three binder facts.}
\label{obs:bound-binder-facts}
The following consequences of the graph definitions will be used
repeatedly.  Suppose a normalized product has profile $K$, domain profile
$I$, open codomain profile $J$, and instantiated codomain profile $L$.
If $K$ and $L$ are active, the positive ascent $L\longrightarrow
J\longrightarrow K$ makes $J$ active.  If also $I\in H$, the retained
negative edge $I\longrightarrow K$ makes $I$ active.  Therefore:
\begin{enumerate}
\item If $K$ is plain, every active frame payload and every active
continuation is plain: an active edge into a plain vertex starts at a
plain vertex, and the same holds for a positive path into a plain vertex.
\item If $K$ and $L$ are free, they have the same sign; a free domain
payload has the opposite sign to $K$.
\item If $K,L\in S$, then $J\in S$, since $J$ lies on a path between two
vertices of the primary strongly connected component $S$.  Hence an
active $I$ is direct.  A free direct vertex is plus, so such a domain is
never free minus.
\end{enumerate}
The last assertion is the reason that universal quantification over
payloads in a product candidate preserves a fixed minus bound.

\begin{lemma}[Uniform finite bounds]\label{obs:bound-statement}
For every fixed finite tree $t$ there are bounds for each of its
$\operatorname{Val}$ and $\operatorname{Read}$ evaluations, denoted below
by $F_t(k)$ when the evaluation is understood, with
$F_t(k)\geq k+1$, satisfying the following assertions.
\begin{enumerate}
\item Every evaluation is defined for every typed environment.
\item In any environment minus-bounded by $k$, a free minus
$\operatorname{Val}$ output belongs to $C_{F_t(k)}$.
\item For two environments minus-bounded by $k$ and agreeing to level
$F_t(k)$, plain and minus $\operatorname{Val}$ outputs are equal,
$\operatorname{Read}$ outputs are equal, and plus
$\operatorname{Val}$ outputs are $\equiv_k$-equivalent.
\end{enumerate}
The output spaces and their indices in the two environments agree:
pointwise convertibility of the images gives the same normalized actual
types, expression profiles, normalized argument keys, and plain type
tags.  Increasing a chosen bound preserves all three assertions.
\end{lemma}
\begin{proof}
We prove the statement together with totality by induction on syntax
size, proving $\operatorname{Val}$ before $\operatorname{Read}$ at a
given size.  All calls to $\operatorname{Read}$ from a product
$\operatorname{Val}$ concern proper children; the only same-size call
is from the remaining $\operatorname{Read}$ clause to
$\operatorname{Val}$.  Thus this is a well-founded order.  A bound for a
node is the maximum of the finitely many bounds described below for its
possible clauses.  Conditional variation of the actual profiles does
not introduce infinitely many cases: the estimates distinguish only
inactive, plain, plus, and minus layers and the finitely many syntactic
evaluation clauses.

\paragraph{Lookup and defaults.}
\label{obs:bound-leaves}
A minus lookup already lies in $C_k$, and a plus lookup agrees under
$\equiv_k$ whenever the environments agree to any level at least $k$.
Plain lookups agree exactly.  Constant-by-ending defaults lie in $C_1$
on the minus side.  Sort defaults and skipped-application defaults
therefore satisfy the assertions with bound $k+1$.

\paragraph{Lambdas at plain and plus layers.}
\label{obs:bound-positive-lambdas}
Write $l$ for the bound of the body evaluation at $k$.  A lambda at a
plain layer is defined using the plain prefix bijection.  Each of its
frames introduces an identical plain parameter, or no parameter, when
two matching frames are compared.  Its active continuation is also
plain, by the first binder fact.  The extended environments remain
minus-bounded by $k$, so body comparison gives equal frame outputs.
Together with the equal default base, the prefix bijection gives equal
lambda values.  This holds simultaneously for all frames and all
continuations; the same $l$ works for every one.

A lambda at a plus layer is a full table.  It is defined on every raw
stack using the gate in the evaluation definition.  For each individual
valid first frame, its minus payload, if present, belongs to some finite
stage of $C_*$.  Together with the finitely many original environment
parameters this gives a finite bound for the extended environment.
Totality of the smaller body tree therefore defines the requested body
value.  This use of totality does not require a single bound for all
possible payloads in the full table.

For comparison under $\equiv_k$, however, only stacks whose minus
payloads all lie in $C_k$ are inspected.  Compare the same such stack in
the two tables.  The gate agrees in the two environments.  If it is
valid, its first frame introduces either the same minus parameter in
$C_k$, the same plain parameter, or none.  Body comparison at $k$ thus
applies.  For a continue frame the body has plus sign and its
$\equiv_k$-agreement gives the same observation on the restricted tail.
For an exit frame the body is plain and agrees exactly.  Invalid stacks
and the empty stack have matching defaults.  This proves
$\equiv_k$-agreement of the plus lambda using the bound $l$.

\paragraph{Lambdas at minus layers.}
\label{obs:bound-minus-lambdas}
Again let $l$ be the body's bound at $k$, so $l\geq k+1\geq 1$.
First define the lambda's behavior on raw stacks by the stated gate and
body clauses.  Every valid first frame introduces a plus parameter, a
plain parameter, or none; no new minus parameter is introduced.  The
body is therefore total in an environment still minus-bounded by $k$.

Consider two original environments agreeing to level $l$, and two
stacks with identical keys, flags, none/plain payloads, and pairwise
$\equiv_l$-equivalent plus payloads.  Their gates agree.  At a valid
first frame the extended environments still agree to level $l$ and
remain minus-bounded by $k$.  If the frame continues, the two body
outputs are the same minus value $c\in C_l$.  Its behavior factors
through $\equiv_{l-1}$, and hence also through $\equiv_l$, on the plus
payloads of the tails.  It therefore gives equal observations on the
two tails.  If the frame exits, body comparison gives equal plain
values.  Defaults and empty observations also agree.

Apply this argument first with the same original environment on both
sides.  It proves uniform invariance of the lambda behavior under
$\equiv_l$ in all plus payloads.  The exact factorization result for
free spaces then represents it by an element of $C_{l+1}$.  Apply the
argument next to different original environments and the same stack on
both sides.  Their represented behaviors agree on every stack; the
behavior representation is injective, so their minus values are equal.
Thus $l+1$ is a valid output and comparison bound for the minus lambda.

\paragraph{Applications.}
\label{obs:bound-applications}
Let $f$ and $n$ be bounds for the function and argument evaluations,
respectively.  Actual normalized argument keys agree in image-compatible
environments.  A function at an inactive layer returns the default.  A
function at a plain layer has only plain active dependencies, by the
first binder fact.  The function and any evaluated argument payload
therefore agree exactly by their induction hypotheses at $k$, and the
plain prefix operation gives equal outputs.

Suppose the function is plus.  Argument comparison at $k$ makes a minus
payload exactly equal in the two environments and puts it in
$C_{n(k)}$.  Set
\[
                 a=\max\{k,n(k)\}.
\]
The original environments are also minus-bounded by $a$.  Function
comparison at $a$ gives $\equiv_a$-agreement, provided the environments
agree to level $f(a)$.  Every prefixed stack needed to compare a
continuing output under $\equiv_k$ has its new minus payload in
$C_{n(k)}\subseteq C_a$ and all tail minus payloads in
$C_k\subseteq C_a$.  The function observations therefore agree on
these stacks.  An exit is a single such frame and gives equal tagged
plain values, hence equal results after projection.  If the payload is
none or plain, the same estimate applies with no use of $C_{n(k)}$.

Suppose instead the function is minus.  Function comparison at $k$
gives the same value in $C_j$, where
\[
                         j=f(k).
\]
Apply argument comparison at $b=\max\{k,j\}$.  If the argument payload
is plus, the two payloads agree under $\equiv_b$, hence under
$\equiv_j$.  Since the function behavior already factors through
$\equiv_{j-1}$, inserting these payloads gives the same observations
on identical arbitrary tails.  A continue slice remains in $C_j$ by
the slicing result, and the two slices are equal by injectivity of the
behavior representation.  An exit gives equal plain values.  Equal
plain payloads and absent payloads are immediate special cases.

All these calls, together with the required output stages, are dominated
by the finite expression
\begin{equation}
  1+\max\bigl\{k,f(k),n(k),
       f(\max\{k,n(k)\}),n(\max\{k,f(k)\})\bigr\}.
  \label{obs:bound-application-formula}
\end{equation}
No monotonicity of $f$ or $n$ is needed: the displayed arguments are
the precise ones used above.  A plus active function can only have a
free plus active continuation, and a minus one only a free minus
active continuation, so the cases prove the required output assertion
as well as comparison.

\paragraph{Product candidates and product overrides.}
\label{obs:bound-products}
Consider a product candidate whose actual expression profile is in $S$.
Use the domain child's $\operatorname{Read}$ comparison at $k$ whenever
that Read is requested.  It gives exactly the same domain candidate in
the two environments, and hence exactly the same set of testers.  If
the domain Read is not requested, the typed strongly normalizing tester
set is already identical because the actual normalized domains agree.

For each common tester $n$, the instantiated actual codomain types
agree.  When their profile lies in $S$, the third binder fact says
that an active domain payload is plain or plus, never minus.  Pair
each payload choice with the identical choice in the other environment.
This preserves the minus bound $k$ and agreement at every level
already required for the original environments.  The codomain child's
$\operatorname{Read}$ comparison at $k$ therefore gives equal codomain
candidates for every tester and every payload choice.  Its bound is
uniform in the raw term $n$ and the payload: both are variable images
or parameters for the same fixed smaller tree.  If the instantiated
profile is outside $S$, the required typed strong-normalization
condition is identical on the two sides.  Thus all the universally
quantified membership conditions defining the product candidate agree,
so the candidate itself is equal.  The maximum of the two child bounds
and $k+1$ suffices.

The product $\operatorname{Val}$ override places this candidate in one
component of the empty-stack $\mathcal B$ observation, leaving the
other components and all nonempty observations at their defaults.
The actual component index is the same in image-compatible environments.
For a minus output such a table is represented already in $C_1$,
because its nonempty observations are constant by ending and its empty
observation contains no payload.  For a plus output its observations
agree exactly, in particular under $\equiv_k$.  At a plain layer the
default continuations and equal base give equality by the prefix
bijection.  This proves the override case with the same child bounds.

\paragraph{The remaining Read clauses.}
\label{obs:bound-read}
The default Read clauses are immediate.  At a nonproduct whose Read
uses its $\operatorname{Val}$ base, use the already proved
same-size $\operatorname{Val}$ bound.  Plain and minus values agree
exactly; plus values agree under $\equiv_k$.  The latter also gives
equal empty observations, even when $k=0$, since the empty stack
belongs to every restriction.  Reading the same actual-expression
component therefore gives the same candidate.  This completes the
simultaneous induction.
\end{proof}

\paragraph{Exact invariance for identical parameters.}
\label{obs:bound-exact-invariance}
As a consequence, image-compatible environments with identical parameters
give identical outputs of every kind.  For plain, minus, and Read
outputs choose a common finite minus bound and use the comparison
statement.  For plus outputs, fix any raw observation stack.  Its
finitely many minus payloads and the finitely many minus parameters of
the original environments lie in some common $C_k$.  The environments
agree to every level because their parameters are identical, so the
comparison statement gives $\equiv_k$-agreement of the plus outputs
and hence equality on this stack.  The stack was arbitrary, proving
equality of the full tables.

The finite-bound proof applies to arbitrary typed trees; it assumes
neither normal syntax nor goodness.  In particular its uniformity is
unaffected by the lengths of reductions needed to normalize varying
raw images.  Nonfunctional axioms and product rules also cause no new
case: the fixed tree stores its selected judgments, while compatibility
aligns the corresponding actual normalized indices.
